\documentclass[preview]{standalone}

\usepackage{amsmath}
\usepackage{amssymb}
\usepackage{stellar}
\usepackage{definitions}

\begin{document}

\id{algebras-over-fields}
\genpage

\section{Algebras over fields}

\begin{snippetdefinition}{algebra-definition}{Algebra over a field}
    Let \(K\) be a \field. A \emph{\(K\)-algebra} is a \set \(A\) that is
    both a \ring and a \vectorspace over \(K\), with the same addition, such
    that
    \[
        \lambda(ab)=(\lambda a)b=a(\lambda b)
    \]
    for all \(\lambda\in K\) and \(a,b\in A\).
\end{snippetdefinition}

\begin{snippet}{algebra-as-subring}
    Let \(A\) be a \(K\)-\algebra. The \function
    \[
        \iota\colon K\fromto A,
        \qquad
        \iota(\lambda)=\lambda1_A,
    \]
    is a unital \ringhomomorphism. It is \injective when \(A\) is
    nontrivial, so in that case \(K\) is identified with the \subring
    \[
        K1_A=\{\lambda1_A\suchthat\lambda\in K\}\subringleq A.
    \]
\end{snippet}

\begin{snippetproposition}{algebra-scalars-lie-in-center}{Scalars lie in the center}
    Let \(A\) be a \(K\)-\algebra. Then
    \[
        \{\lambda1_A\suchthat\lambda\in K\}\subseteq Z(A).
    \]
\end{snippetproposition}

\begin{snippetproof}{algebra-scalars-lie-in-center-proof}{algebra-scalars-lie-in-center}{Scalars lie in the center}
    For every \(a\in A\) and \(\lambda\in K\),
    \[
        (\lambda1_A)a=\lambda(1_Aa)=\lambda a
        =\lambda(a1_A)=a(\lambda1_A).
    \]
    Hence \(\lambda1_A\in Z(A)\).
\end{snippetproof}

\section{Examples}

\begin{snippetexample}{algebra-example1}{Polynomial algebras}
    Let \(K\) be a \field. The polynomial \ring \(K[x]\) is a
    \(K\)-\algebra. The copy of \(K\) inside \(K[x]\) consists of the
    constant polynomials.
\end{snippetexample}

\begin{snippetexample}{free-noncommutative-algebra-example}{A free noncommutative algebra}
    Let \(K\) be a \field. The \set \(K\langle X,Y\rangle\) of
    noncommutative polynomials in \(X\) and \(Y\) is a \(K\)-\algebra.
    Its elements are finite \(K\)-linear combinations of words in \(X\) and
    \(Y\), for example
    \[
        \lambda_1XY^2X+\lambda_2YXX+\lambda_3X+\lambda_4.
    \]
    In general \(XY\neq YX\), and
    \[
        Z(K\langle X,Y\rangle)=K.
    \]
\end{snippetexample}

\begin{snippetexample}{matrix-algebra-example}{The matrix algebra}
    Let \(K\) be a \field. The matrix \ring \(\matrices_n(K)\) is a
    \(K\)-\algebra. Its scalar matrices form its center:
    \[
        Z(\matrices_n(K))
        =\{\lambda I_n\suchthat\lambda\in K\}.
    \]
    If \(\Delta_n(K)\) denotes the \subring of diagonal matrices, then
    \[
        C_{\matrices_n(K)}(\Delta_n(K))=\Delta_n(K).
    \]
\end{snippetexample}

\begin{snippetdefinition}{group-algebra-definition}{Group algebra}
    Let \(K\) be a \field and let \(G\) be a finite \group. The
    \emph{group algebra of \(G\) over \(K\)} is the \(K\)-\vectorspace
    \[
        K[G]
        =\left\{\sum_{g\in G}\lambda_g g
        \suchthat \lambda_g\in K\right\},
    \]
    with componentwise addition and multiplication obtained by bilinearly
    extending the multiplication of \(G\):
    \[
        \left(\sum_{g\in G}\lambda_g g\right)
        \left(\sum_{h\in G}\mu_h h\right)
        =\sum_{g,h\in G}\lambda_g\mu_h(gh).
    \]
    Its multiplicative identity is \(1_Ke_G\).
\end{snippetdefinition}

\begin{snippetexample}{cyclic-group-order-two-algebra-example}{The group algebra of the cyclic group of order two}
    Let \(C_2=\{1,\sigma\}\), where \(\sigma^2=1\). Then
    \[
        K[C_2]=\{\lambda+\mu\sigma\suchthat\lambda,\mu\in K\}.
    \]
    For \(\alpha_i=\lambda_i+\mu_i\sigma\),
    \[
        \alpha_1\alpha_2
        =(\lambda_1\lambda_2+\mu_1\mu_2)
        +(\lambda_1\mu_2+\lambda_2\mu_1)\sigma.
    \]
\end{snippetexample}

\begin{snippetexample}{quaternions-real-algebra-example}{The quaternions as a real algebra}
    The quaternion division ring
    \[
        \mathbb{H}
        =\operatorname{span}_{\realnumbers}\{1,i,j,k\}
    \]
    is an \(\realnumbers\)-\algebra, and
    \[
        Z(\mathbb{H})
        =\operatorname{span}_{\realnumbers}\{1\}
        \cong\realnumbers.
    \]
    The real subspaces
    \[
        A_1=\operatorname{span}_{\realnumbers}\{1,i\},
        \quad
        A_2=\operatorname{span}_{\realnumbers}\{1,j\},
        \quad
        A_3=\operatorname{span}_{\realnumbers}\{1,k\}
    \]
    are real subalgebras of \(\mathbb{H}\), each isomorphic to
    \(\complexnumbers\) as an \(\realnumbers\)-\algebra. Moreover,
    \[
        C_{\mathbb{H}}(A_1)=A_1.
    \]
\end{snippetexample}

\end{document}
